% /chapters/3-deep-computations/05-randomized.tex

\section{Monte Carlo computability}\label{sec:measure-theoretic-NIP}

In this section, we replace deterministic computability by probabilistic (`Monte Carlo') computability.
We do not assume that $\mathcal{P}$ is countable.
(The countability assumption on $\mathcal{P}$ played a crucial role in the proof of Theorem~\ref{thm:Generalized-BFT}, as it makes $\mathbb{R}^\mathcal{P}$ a Polish space.)
The main results of the section are Theorem~\ref{thm:NIP-Monte-Carlo}, which connects NIP and Monte Carlo computability, and~\ref{cor:Talagrand-stab-Monte-Carlo}, which connects Talagrand stability and Monte Carlo computability.

Fundamental in this section is a measure-theoretic version of Theorem~\ref{thm:Generalized-BFT}, namely, Theorem~\ref{thm:BFT-2F}.

\subsection{NIP and Monte Carlo computability of deep computations}

By Fact~\ref{baire}, for perfectly normal $X$, a function $f$ is in $\mathrm{B}_1(X,Y)$ if and only if $f^{-1}[U]$ is an $F_\sigma$ subset of $X$ for every open $U\subseteq Y$.
This motivates the following definition.

\begin{definition}\label{def:universal-measurability}
    Given a Hausdorff space $X$ and a measurable space $(Y,\Sigma)$, we say that $f\colon X\to Y$ is \emph{universally measurable} (with respect to $\Sigma$) if $f^{-1}(E)$ is $\mu$-measurable for every Radon measure $\mu$ on $X$ and every $E\in\Sigma$.
    When $Y=\mathbb{R}$, we will always take $\Sigma=\mathcal{B}(\mathbb{R})$, the Borel $\sigma$-algebra of $\mathbb{R}$.
\end{definition}

If $X$ is a compact (Hausdorff) space, then Radon measures on $X$ are finite, and hence, through renormalization, we can regard them as probability measures.
We thus have the following remark:

\begin{remark}\label{rem:probability}
If $X$ is compact, then a function $f\colon X\to\mathbb{R}$ is universally measurable if and only if $f^{-1}(U)$ is $\mu$-measurable for every Radon probability measure $\mu$ on $X$ and every open set $U\subseteq\mathbb{R}$.
\end{remark}

Intuitively, a function is universally measurable if it is ``measurable no matter which reasonable way you try to measure things on its domain".
The concept of universal measurability emerged from work of Kallianpur and Sazonov, in the late 1950's and 1960s --- with later developments by Blackwell, Darst and others --- building on earlier ideas of Gnedenko and Kolmogorov from the 1950s.
See~\cite[Chapters 1 and 2]{Pap:2002}.

\begin{notation}
    Following \cite{BFT_1978_PCompactBaire}, the collection of all universally measurable real-valued functions on $X$ will be denoted by $M_r(X)$.
    Given a fixed Radon measure $\mu$ on $X$, the collection of all $\mu$-measurable real-valued functions on $X$ will be denoted by $\mathscr{M}^0(X,\mu)$.
\end{notation}

In the context of deep computations, we are interested in transition maps of a state space $L\subseteq \mathbb{R}^\mathcal{P}$ into itself.
In the product space $\mathbb{R}^\mathcal{P}$, we have two natural $\sigma$-algebras: 
the Borel $\sigma$-algebra, i.e., the $\sigma$-algebra generated by open sets in $\mathbb{R}^\mathcal{P}$, and the cylinder $\sigma$-algebra (i.e., the $\sigma$-algebra generated by the sub-basic open sets in $\mathbb{R}^\mathcal{P}$, that is, the set $\pi_i^{-1}(U)$ with $U\subseteq \mathbb{R}$ open).
In general, the cylinder $\sigma$-algebra is strictly smaller; however, when $\mathcal{P}$ is countable, both $\sigma$-algebras coincide.
We will use the cylinder $\sigma$-algebra to define universally measurable maps $f\colon\mathbb{R}^\mathcal{P}\to\mathbb{R}^\mathcal{P}$.
The reason for this choice is the following characterization:

\begin{proposition}\label{prop:coordinated_univ_meas}
    Let $X$ be a Hausdorff space and $Y=\prod_{i\in I}Y_i$ be any product of measurable spaces $(Y_i,\Sigma_i)$ for $i\in I$.
    Let $\Sigma_Y$ be the cylinder $\sigma$-algebra generated by the measurable spaces $(Y_i,\Sigma_i)$.
    The following are equivalent for $f\colon X\to Y$:
    \begin{enumerate}[(i)]
        \item $f\colon X\to Y$ is universally measurable (with respect to $\Sigma_Y$).
        \item $\pi_i\circ f\colon X\to Y_i$ is universally measurable (with respect to $\Sigma_i$) for all $i\in I$.
    \end{enumerate}
\end{proposition}

\begin{proof}
    (i)$\Rightarrow$(ii) is clear since the projection maps $\pi_i$ are measurable and the composition of measurable functions is measurable.
    To prove (ii)$\Rightarrow$(i), note that if $C=\prod_{i\in I}C_i$ is a measurable cylinder and $J$ is the finite subset of $I$ such that $C_i\neq Y_i$ for $i\in J$, then, $f^{-1}(C)=\bigcap_{i\in J}(\pi_i\circ f)^{-1}(C_i)$ is universally measurable by assumption.
\end{proof}

The preceding proposition says that we can check measurability of a transition just by checking measurability feature by feature.
This is analogous to the Baire class 1 setting (see Proposition~\ref{prop:Baire-identifications}).

We saw in section \ref{sec:prelim} that, when the set $\mathcal{P}$ of predicates is countable, PAC-learning (or NIP) corresponds to relative compactness in the space of Baire class 1 functions (see Theorems~\ref{thm:BFT}, \ref{thm:BFT2}, \ref{thm:Generalized-BFT}).
The following result (which does not assume countability of $\mathcal{P}$) gives an analogous characterization of the NIP in terms of universal measurability:

\begin{theorem}
[Bourgain-Fremlin-Talagrand, Theorem 2F in \cite{BFT_1978_PCompactBaire}]\label{thm:BFT-2F}
    Let $X$ be a Hausdorff space and $A\subseteq \mathrm{C}(X)$ be pointwise bounded.
    The following are equivalent:
    \begin{enumerate}[(i)]
        \item $\overline{A}\subseteq M_r(X)$.
        \item For every compact $K\subseteq X$, $A\restriction_K$ satisfies the NIP.
        \item For every Radon measure $\mu$ on $X$, $A$ is relatively countably compact in $\mathscr{M}^0(X,\mu)$, i.e., every countable subset of $A$ has a limit point in $\mathscr{M}^0(X,\mu)$.
    \end{enumerate}
\end{theorem}

This result allows us to formalize the concept of a deep computation being \emph{Monte Carlo computable}, which we define below (Definition~\ref{def:universally-essentially-computable}).
To motivate the definition, let us first recall two facts:

\begin{enumerate}
    \item Littlewoood's second principle states that every Lebesgue measurable function is ``nearly continuous.''
The formal statement of this, which is Luzin's theorem, is that if $(X,\Sigma ,\mu )$ a Radon measure space and $Y$ is a second-countable topological space (e.g., $Y=\mathbb{R}^\mathcal{P}$ with $\mathcal{P}$ countable) equipped with the Borel $\sigma$-algebra, then any given $f:X\to Y$ is measurable if and only if for every $E\in\Sigma$ and every $\varepsilon>0$ there exists a closed $F\subseteq E$ such that the restriction $f\restriction_F$ is continuous and $\mu(E\backslash F) < \varepsilon$.

    \item Computability of deep computations is characterized in terms of continuous extensibility of computations.
    This is at the core of~\cite{alva2024approximability}.
\end{enumerate}

These two facts motivate the following definition:

\begin{definition}\label{def:universally-essentially-computable}
    Let $(L,\mathcal{P},\Gamma)$ be a CCS.
    We say that a transition $f\colon L\to L$ is \emph{universally Monte Carlo computable} if and only if there exists an extension $\tilde f:\colon\mathcal{L}_{\sh}\to \mathcal{L}_{\sh}$ of $f$ satisfying the following property: for every sizer $r_\bullet$ there is a sizer $s_\bullet$ such that the restriction $\tilde f\restriction_{\mathcal{L}[r_{\bullet}]}\colon\mathcal{L}[r_\bullet]\to\mathcal{L}[s_\bullet]$ is universally measurable, i.e., $\pi_P\circ \tilde f\restriction_{\mathcal{L}[r_{\bullet}]}\colon \mathcal{L}[r_\bullet]\to [-s_P,s_P]$ is $\mu$-measurable for every Radon probability measure $\mu$ on $\mathcal{L}[r_\bullet]$ and $P\in\mathcal{P}$.
\end{definition}

Note that, since $\mathcal{L}_{\sh}$ is compact, by remark~\ref{rem:probability}, in the context of Monte Carlo computability, we restrict our attention to Radon probability measures instead of general Radon measures.

\begin{theorem}\label{thm:NIP-Monte-Carlo}
    Let $(L,\mathcal{P},\Gamma)$ be a CCS satisfying the Extensibility Axiom.
    Let $R$ be an exhaustive collection of sizers.
    Let $\Delta\subseteq\Gamma$ be $R$-confined.
    If $\pi_P\circ\Delta\restriction_{L[r_\bullet]}$ satisfies the NIP for all $P\in\mathcal{P}$ and all $r_{\bullet}\in R$, then every deep computation in $\Delta$ is universally Monte Carlo computable.
\end{theorem}

\begin{proof}
    Fix $P\in\mathcal{P}$ and $r_\bullet\in R$.
    By the Extensibility Axiom, $\pi_P\circ\tilde\Delta\restriction_{\mathcal{L}[r_\bullet]}$ is a set of pointwise bounded continuous functions on the compact set $\mathcal{L}[r_\bullet]$.
    Since $\pi_P\circ\Delta\restriction_{L[r_\bullet]}$ has the NIP, so does $\pi_P\circ\tilde\Delta\restriction_{L[r_\bullet]}=\pi_P\circ\Delta\restriction_{\mathcal{L}[r_\bullet]}$, by Lemma \ref{lem:NIP-and-closure}.
    Let $f\in\overline{\Delta}$ be a deep computation;
     write $f$ as an ultralimit of computations in $\Delta$, say $f=\mathcal{U}\lim_i\gamma_i$ , and
    define an extension $\tilde{f}\colon\mathcal{L}_{\sh}\to \mathcal{L}_{\sh}$ of $f$ naturally by $\tilde{f}:=\mathcal{U}\lim_i\tilde{\gamma_i}$.
    Since, by hypothesis, $\Delta$ is $R$-confined, we have $f\colon L[r_\bullet]\to L[r_\bullet]$ and hence, $\tilde{f}\colon\mathcal{L}[r_\bullet]\to \mathcal{L}[r_\bullet]$.
    Now, by Theorem \ref{thm:BFT-2F}, we have $\overline{\pi_P\circ\tilde\Delta\restriction_{\mathcal{L}[r_\bullet]}}\subseteq M_r(\mathcal{L}[r_\bullet])$; therefore, $\pi_P \circ \tilde{f}\restriction_{\mathcal{L}[r_\bullet]}\in \overline{\pi_P\circ\tilde\Delta\restriction_{\mathcal{L}[r_\bullet]}}\subseteq M_r(\mathcal{L}[r_\bullet])$.
\end{proof}

\begin{question}
    Under the same assumptions of the preceding theorem, suppose that every deep computation of $\Delta$ is universally Monte Carlo computable.
    Must $\pi_P\circ\Delta\restriction_{L[r_\bullet]}$ have the NIP for all $P\in\mathcal{P}$ and all $r_{\bullet}\in R$?
\end{question}

\subsection{Talagrand stability and Monte Carlo computability of deep computations}

There is another notion closely related to NIP, introduced by Talagrand in \cite{talagrand1984pettis} while studying Pettis integration.
Suppose that $X$ is a compact Hausdorff space and $A\subseteq \mathbb{R}^X$.
Let $\mu$ be a Radon probability measure on $X$.
Given a $\mu$-measurable set $E\subseteq X$, a positive integer $k$ and real numbers $a<b$, let

\[
	D_k(A,E,a,b)=\bigcup_{f\in A}\{x\in E^{2k}:f(x_{2i})\leq a, \hspace{1mm} f(x_{2i+1})\geq b \hspace{1mm}{\text{ for all }i<k} \}.
\]

\begin{definition}\label{def:Talagrand-stable}
    We say that $A$ is \emph{Talagrand $\mu$-stable} if and only if for every $\mu$-measurable set $E\subseteq X$ of positive measure and for every $a<b$ there is $k\geq 1$ such that
    \[
    (\mu^{2k})^*(D_k(A,E,a,b))<(\mu(E))^{2k},
    \]
    where $\mu^*$ denotes the outer measure.
\end{definition}

Note that the reference to outer measure in Definition~\ref{def:Talagrand-stable} is necessary since $D_k(A,E,a,b)$ need not be measurable.
The inequality certainly holds when $A$ is a countable set of continuous (or $\mu$-measurable) functions.

\begin{remark}\label{rem:Talgrand-stability}
    Let $X$, $\mu$, $A\subseteq \mathbb{R}^X$, and $\subseteq X$ be as above. Then,
    \begin{enumerate}
        \item If $A\subseteq B\subseteq \mathbb{R}^X$, then \[D_k(A,E,a,b)\subseteq D_k(B,E,a,b).\]
        \item If $a<a'<b'<b$, then \[D_k(\overline{A},E,a,b)\subseteq D_k(A,E,a',b').\]

[\emph{Proof}: If $x\in D_k(\overline{A},E,a,b)$, then there is $f\in\overline{A}$ such that $f(x_{2i})\leq a<a'$ and $f(x_{2i+1})\geq b>b'$ for all $i<k$.
By definition of pointwise convergence topology, there exists $g\in A$ such that $g(x_{2i})<a'$ and $g(x_{2i+1})>b'$ for all $i<k$.
Hence, $x\in D_k(A,E,a',b')$.]

    \item Any subset of a $\mu$-stable set is $\mu$-stable. (This is immediate from (1).)

    \item If $A$ is stable, then $\overline{A}$ is $\mu$-stable. (This is immediate from (2).)
\end{enumerate}
\end{remark}

The main result of this section is that deep computations that are limit points of a Talagrand stable set of computations are Monte Carlo computable; this is Corollary~\ref{cor:Talagrand-stab-Monte-Carlo} below.
To prove it, we first show that the closure of a Talagrand $\mu$-stable set is $\mu$-measurable (lemma~\ref{lem:talagrand-stab-relcpct});
but first, we need the following useful characterization of measurable functions (compare with Fact \ref{baire}):

\begin{fact}[Lemma 413G in \cite{fremlin2003vol4}]\label{fact:meas-fns-charn}
    Suppose that $(X,\Sigma,\mu)$ is a measure space and let $\mathcal{K}\subseteq\Sigma$ be a collection of measurable sets satisfying the following conditions:
    \begin{itemize}
        \item $(X,\Sigma,\mu)$ is complete, i.e., for all $E\in\Sigma$ with $\mu(E)=0$ and $F\subseteq E$ we have $F\in\Sigma$.
        \item $(X,\Sigma,\mu)$ is semi-finite, i.e., for all $E\in\Sigma$ with $\mu(E)=\infty$ there exists $F\subseteq E$ such that $F\in\Sigma$ and $0<\mu(F)<\infty$.
        \item $(X,\Sigma,\mu)$ is saturated, i.e., $E\in\Sigma$ if and only if $E\cap F\in\Sigma$ for all $F\in\Sigma$ with $\mu(F)<\infty$.
        \item $(X,\Sigma,\mu)$ is inner regular with respect to $\mathcal{K}$, i.e., for all $E\in\Sigma$
        \[
	\mu(E)=\sup\{\mu(K):K\in\mathcal{K} \text{ and } K\subseteq E\}.
        \]
    \end{itemize}
    (In particular, if $X$ is compact Hausdorff, $\mathcal{K}$ is the collection of compact subsets of $X$, $\mu$ is a Radon probability measure on $X$, $\Sigma$ is the completion of the Borel $\sigma$-algebra by $\mu$, and all these conditions hold).
    Then, $f\colon X\to\mathbb{R}$ is measurable if and only if for every $K\in\mathcal{K}$ with $0<\mu(K)<\infty$ and $a<b$, either $\mu^*(P)<\mu(K)$ or $\mu^*(Q)<\mu(K)$, where $P=\{x\in K:f(x)\leq a\}$ and $Q=\{x\in K:f(x)\geq b\}$.
\end{fact}

The following technical lemma will be instrumental for proving Proposition~\ref{prop:Talagrand-implies-NIP}, which, in turn, will yield the main result of the subsection, namely Corollary~\ref{cor:Talagrand-stab-Monte-Carlo}

\begin{lemma}\label{lem:talagrand-stab-relcpct}
    If $A$ is Talagrand $\mu$-stable, then $\overline{A}$ is also Talagrand $\mu$-stable and $\overline{A}\subseteq\mathscr{M}^0(X,\mu)$.
\end{lemma}

\begin{proof}
    If $A$ is Talagrand $\mu$-stable, so is $\overline{A}$, by Remarks~\ref{rem:Talgrand-stability}.
    Suppose that $A$ is Talagrand $\mu$-stable to show that $\overline{A}\subseteq \mathscr{M}^0(X,\mu)$.
    Assume, for the sake of contradiction that there exists $f\in\overline{A}$ such that $f\notin \mathscr{M}^0(X,\mu)$.
    By Fact \ref{fact:meas-fns-charn}, there exists a $\mu$-measurable set $E$ of positive measure and $a<b$ such that $\mu^*(P)=\mu^*(Q)=\mu(E)$ where $P=\{x\in E: f(x)\leq a\}$ and $Q=\{x\in E: f(x)\geq b\}$.
    Then, for any $k\geq 1$: $(P\times Q)^k\subseteq D_k(\{f\},E,a,b)$, so $(\mu^{2k})^*(D_k(\{f\},E,a,b))=(\mu^*(P)\mu^*(Q))^k=(\mu(E))^{2k}$.
    Thus, the set $\{f\}$ is not $\mu$-stable. 
    However, $\{f\}\subseteq\overline{A}$ and we know (Remarks~\ref{rem:Talgrand-stability}) that a subset of a $\mu$-stable set must be $\mu$-stable,
    so we have a contradiction.
\end{proof}

\begin{definition}
We say that $A$ is \emph{universally Talagrand stable} if $A$ is Talagrand $\mu$-stable for every Radon probability measure $\mu$ on $X$.
\end{definition}

We first observe that universal Talagrand stability implies NIP:

\begin{proposition}\label{prop:Talagrand-implies-NIP}
    Let $X$ be a compact Hausdorff space and $A\subseteq \mathrm{C}(X)$ be pointwise bounded.
    If $A$ is universally Talagrand stable, then $A$ satisfies the NIP.
\end{proposition}

\begin{proof}
    Suppose that $A$ is universally Talagrand stable.
    By Theorem \ref{thm:BFT-2F}, it suffices to show that $A$ is relatively countably compact in $\mathscr{M}^0(X,\mu)$ for every Radon probability measure $\mu$ on $X$.
    Since, for any such $\mu$, $A$ is Talagrand $\mu$-stable, Lemma \ref{lem:talagrand-stab-relcpct} gives us $\overline{A}\subseteq\mathscr{M}^0(X,\mu)$.
    In particular, $A$ is relatively countably compact in $\mathscr{M}^0(X,\mu)$.
\end{proof}

\begin{corollary}
\label{cor:Talagrand-stab-Monte-Carlo}
    Let $(L,\mathcal{P},\Gamma)$ be a CCS satisfying the Extensibility Axiom.
If $\pi_P\circ\Delta\restriction_{L[r_\bullet]}$ is universally Talagrand stable for all $P\in\mathcal{P}$ and all sizers $r_{\bullet}$, then every deep computation is universally Monte Carlo computable.
\end{corollary}

\begin{proof}
    This follows directly from Proposition \ref{prop:Talagrand-implies-NIP} and Theorem~\ref{thm:NIP-Monte-Carlo}.
\end{proof}

In the context of deep computations, we have identified two ways to obtain Monte Carlo computability, namely, NIP/PAC and Talagrand stability.
It is natural to ask whether these two notions are equivalent. 
The following results show that, even in the simple case of countably many computations, this question is sensitive to the set-theoretic axioms. 
On the one hand, it is consistent (with respect to the standard ZFC axioms of set theory) that these two classes are the same:

\begin{theorem}[Talagrand, Theorem 9-3-1(a) in \cite{talagrand1984pettis}]\label{thm:NIP-implies-Talagrand-stab}
    Let $X$ be a compact Hausdorff space and $A\subseteq M_r(X)$ be countable and pointwise bounded.
    Assume that $[0,1]$ is not the union of $<\mathfrak{c}$ closed measure zero sets.
If $A$ satisfies the NIP, then $A$ is universally Talagrand stable.
\end{theorem}

(The assumption that $[0,1]$ is not the union of $<\mathfrak{c}$ closed measure zero sets is a consequence of, for example, the Continuum Hypothesis.)

On the other hand, the opposite is also consistent for the particular case of the Lebesgue measure:

\begin{theorem}[Fremlin, Shelah, \cite{fremlin1993pointwise}]
    It is consistent with the usual axioms of set theory that there exists a countable pointwise bounded set of Lebesgue measurable functions with the NIP which is not Talagrand stable with respect to Lebesgue measure.
\end{theorem}

Notice that the preceding two results apply to sets of measurable functions, a class of functions larger than the class of continuous functions.
However, by the Extensibility Axiom, finitary computations are continuous, i.e., if $A$ is a set of computations, then $A\subseteq \Cp(X)$.